\documentclass[]{exercise}

\usepackage{dirtree}
\usepackage{hyperref}
\usepackage{verbatim}
\usepackage{graphicx}
\usepackage{pgffor}

\setgroup{Gruppe 01}

\addstudent[770220]{Luke Berger}
\addstudent[773793]{Benedikt Schröder}
\addstudent[777951]{Alexander Roskamp}

\settitle{Anforderungsdokument - \\ Crazy Machines (MadBalls)}

\begin{document}
    \section{Einleitung:}
        \begin{enumerate}[I.]
            \item \textbf{Hintergrund}\\
                Im Rahmen des Praktikums soll ein reduzierter Klon des Puzzlespiels Crazy Machines (2005) entwickelt werden. Ziel ist es, eine Physik basierte Umgebung bereitzustellen, in der KI-Agenten ihre Fähigkeiten im Bereich Physical Reasoning evaluieren können.
            \item \textbf{Zweck dieses Dokuments}\\
                 Dieses Dokument fasst alle funktionalen Anforderungen zusammen, definiert den Projektumfang und stellt die vorgeschlagene Verzeichnisstruktur für das Repository dar.
        \end{enumerate}

    \section{Zielsetzung:}
        \begin{enumerate}[I.]
            \item \textbf{Primäres Ziel}\\
                Bereitstellung eines stabilen Prototyps mit Grundfunktionen: Physiksimulation, Editiermodus, Puzzle- und Sandbox-Modus.
            \item \textbf{Sekundäres Ziel}\\
                Erweiterung des Prototyps durch optionale Features (Menüs, erweiterter Level-Editor, Grafikalternativen) zur Steigerung von Bedienkomfort, Wiederverwendbarkeit und visueller Identität.
        \end{enumerate}

\begin{comment}
\section{Rahmenbedingungen:}
        \begin{enumerate} [I.]
            \item 
        \end{enumerate}
    I. Teamgröße
        Verfügbar sind 3- oder 4-köpfige Gruppen.

    II. Release-Plan
        Nach Release II werden drei weitere Anforderungen bekannt gegeben.
        • 3er-Gruppen setzen eine der drei nachgereichten Anforderungen um.
        • 4er-Gruppen setzen zwei der drei nachgereichten Anforderungen um.
    III. Bewertung
        Nicht angeforderte oder offensichtlich überflüssige Features können zu Punktabzug führen.
\end{comment}    


    \section{Anforderungen:}
        Folgende Featuers müssen vollständig umgesetzt werden:
        
    \subsection{Pflichtfeatures:}
        
        \begin{enumerate} [I.]
            \item \textbf{Low-Fidelity-Klon}\\
                Reproduktion der Kernspielmechanik gemäß Kick-off-Video.
            \item \textbf{Physiksimulation und Editiermodus}\\
                Die zwei Hauptspielmodi, in denen das Spiel sich befinden kann. Dabei soll es möglich sein:
                \begin{itemize}
                    \item zwischen Editiermodus (Objekte platzieren/entfernen) und Simulationsmodus (Play/Stop) zu wechseln.
                    \item mit Buttons zu navigieren: Start/Abbruch der Simulation, Zurücksetzen in den Editiermodus.
                \end{itemize}
            \item \textbf{Physikobjekte}\\
                Folgende Arten an Objekten mit pysiksimulation sollten existieren:
                \begin{itemize}
                    \item Statische und dynamische Objekte.
                    \item Spezialobjekte (Antriebe, Sensoren) gemäß Kick-off-Folien.
                \end{itemize}
            \item \textbf{Level-System und Puzzle-Modus}\\
                Der Puzzelmodus soll mit folgenden Features erstellt werden:
                \begin{itemize}
                    \item Levels als externe, menschenlesbare Dateien (JSON) speichern
                    \item Import/Export und Verteilung via Dateien
                    \item Automatische Level-Progression bei erfolgreichem Abschluss
                    \item Siegesbedingungen, Einschränkungszonen, Inventarsystem für verfügbare Bauteile
                \end{itemize}
            \item \textbf{Sandboxmodus}:\\
                Zum erleichtern der Erstellung des Puzzle-Modus wollen wir einen Spielplatz/Sandkiste (Sandbox) erstellen, die alle Objekte platzierbar macht und es ermöglicht eigene Level zu bauen.
                \begin{itemize}
                    \item Leere Welt zum freien Bauen
                    \item Import/Export von Levels im Sandbox-Kontext
                    \item Optionale Puzzle-Features (Inventar, Einschränkungszonen) in der Sandbox
                \end{itemize}
            \item \textbf{Grafik}\\
                In der ersten Version genügen Primitive Formen als Objekte, wie Rechtecke oder Kreise, die in verschiedenen Farben angezeigt werden und mit einer Framerate $\geq 30$ in echtzeit gerändert werden.
            \item \textbf{Persistenz}\\
                Es soll dem Spieler per Import/Export von Level Dateien möglich sein, ein Level zu speichern und zu laden und der Spielzustand soll gespeichert werden.
        \end{enumerate}

    \subsection{Wahlfeatures:}
        Zusätzlich zu den Pflichtfeatures sollen drei Wahlfeatures aus 5 gewählt werden. Dabei wählten wir Features \rom{1} Richtige Menüs, \rom{4} Level-Editor+ und \rom{5} komplexere Graphiken.
        \begin{enumerate} [I.]
            \item \textbf{Menüs \& Steuerung}\\
                Hinzufügen von Menü Ansichten, die das Spielerlebnis verbessern:
                \begin{itemize}
                    \item Titelbildschirm mit Start, Puzzle/Sandbox-Auswahl, Optionen (potentiell Audioeinstellungen, Texture-Pack auswählen)
                    \item In-Game-Menü (Pause, Speichern, Beenden)
                \end{itemize}
            \item \textbf{Level-Editor+}\\
                Wir erweiternden Editor, um die Art mit den Objekten im Editor zu interagieren zu verbessern.
                \begin{itemize}
                    \item Drag \& Drop um die Objekte zu platzieren und zu verschieben
                    \item Verschieben und Rotieren von bereits platzierten Objekten
                    \item Eine Undo/Redo-Funktion hinzufügen
                \end{itemize}
            \item \textbf{Komplexere Grafiken}\\
                Wir wollen das Spiel anschaulicher darstellen mit erweiterten Grafikoptionen:
                \begin{itemize}
                    \item Objekte als PNG-Texturen oder Sprites
                    \item Austauschbare Texture-Packs
                    \item Hintergrundgraphiken
                \end{itemize}
            \end{enumerate}
            
        Im weiteren Verlauf wird es weitere Features geben, für die wir auf Wunsch eingehen können müssen. Deshalb sollte weites gehend die Spielstruktur anpassbar sein.

\newpage

    \section{Storycards:}
        Um die gegebenen Anforderungen aus Nutzersicht besser zu verstehen, nutzen wir Storycards (Vorderseite | Rückseite):\\

        \pdfximage{Sources/MadBalls_StoryCards.pdf}
        \newcount\pdfseiten
        \pdfseiten=\pdflastximagepages
        
        \foreach \x in {1,...,\the\pdfseiten} {
        \begin{minipage}{0.5\textwidth}
            \includegraphics[page=\x, width=\textwidth]{Sources/MadBalls_StoryCards.pdf}
        \end{minipage}}

\newpage

    \section{Anwendungsfalldiagramm:}
        Aus diesen Storycards und den oben gegebenen Features erstellten wir nun ein Anwendungsfalldiagramm:
        \begin{figure}[h]
            \centering
            \includegraphics[width=\linewidth]{Sources/Use-Case.png}
            \caption{Anwendungsfalldiagramm}
            \label{fig:Use-Case}
        \end{figure}
        
\newpage

    \section{Verzeichnisstruktur:}
        Die bisher aus den Anforderungen hervorgegangene Struktur haben wir in einem Ordnerbaum dargestellt:
        \dirtree{%
           .1 project\_root/.
            .2 src/.
              .3 core/.
                .4 physics/.
                .4 simulation/.
              .3 editor/.
                .4 tools/.
                .4 ui/.
              .3 game/.
                .4 level/.
                .4 inventory/.
              .3 ui/.
                .4 menu/.
                .4 widgets/.
            .2 assets/.
              .3 textures/.
              .3 backgrounds/.
              .3 fonts/.
            .2 levels/.
            .2 docs/.
              .3 requirements.pdf.
              .3 sanity-checks/.
              .3 design/.
            .2 scripts/.
              .3 build.sh.
            .2 tests/.
              .3 core/.
              .3 editor/.
              .3 game/.
            .2 .gitignore. 
            .2 README.md. 
            .2 pom.xml.   
        }
    

    \section{Glossar:}
    \begin{itemize}
        \item Low-Fidelity: \textit{Funktionale, einfache Darstellung ohne aufwändige Grafiken.}
        \item Persistenz: \textit{Speichern und Laden von Level-Zuständen.}
        \item Sandbox-Modus: \textit{Freies Bauen ohne vorgegebene Puzzle-Ziele.}
        \item Inventarsystem: \textit{Auswahlmenü für verfügbare Bauteile je Level.}
    \end{itemize}
\end{document}
